\chapter{三角形与角}\label{ch:g7:triangles}

能否画一个三边长分别为 $10$、$3$、$2$ 的三角形？能否画一个
三个内角分别为 $80^\circ$、$70^\circ$、$50^\circ$ 的三角形？本章用
平面几何中两个最有用的事实来回答这两个问题：三角形不等式
与 $180^\circ$ 内角和定理。此外，本章还介绍
相交直线与平行线所涉及的角的词汇。

\section{相交直线周围的角}

\begin{definition}[成对的角]\label{def:g7:triangles:pairs}
\begin{itemize}
  \item 两角的度数之和为 $90^\circ$ 时，称为\emph{互余}；若两角的度数之和为
        两角的度数之和为
        $180^\circ$ 时，称为\emph{互补}。
  \item 两条直线相交时，交点两侧相对的
        角称为\emph{对顶
        角}\index{对顶角}。
  \item 一条直线与另外两条直线相交时，
        在每个交点处位置相同的角称为
        \emph{同位角}；位于两直线之间、
        截线两侧的角称为\emph{内错
        角}\index{内错角}。
\end{itemize}
\end{definition}

\begin{theorem}[角的相等关系]\label{thm:g7:triangles:equalities}
\begin{enumerate}
  \item 对顶角相等。
  \item 若两直线平行，则同位角相等，
        内错角也相等。反之，若同位角
        （或内错角）相等，则两直线平行。
\end{enumerate}
\end{theorem}

\begin{proof}[1 的证明]
位于直线同一侧的两个角 $\widehat{1}$ 与 $\widehat{2}$ 构成一个平角，
即 $\widehat 1 + \widehat 2 = 180^\circ$。
$\widehat 3$ 是 $\widehat 1$ 的对顶角，同样满足
$\widehat 2 + \widehat 3 = 180^\circ$（理由相同）。于是
$\widehat 1 = 180^\circ - \widehat 2 = \widehat 3$。第 2 条
在本阶段不加证明地承认。
\end{proof}

\begin{omfigure}
\begin{tikzpicture}[scale=1.0]
  \draw[thick] (-2.0,-0.9) -- (2.0,0.9);
  \draw[thick] (-2.0,0.9) -- (2.0,-0.9);
  \draw[omDef, ->] (0.75,0) arc (0:24.2:0.75);
  \draw[omDef, ->] (-0.75,0) arc (180:204.2:0.75);
  \node[omDef, font=\small] at (1.15,0.28) {$\widehat 1$};
  \node[omDef, font=\small] at (-1.15,-0.28) {$\widehat 3$};
  \draw[omThm, ->] (0.55,0.25) arc (24.2:155.8:0.6);
  \node[omThm, font=\small] at (0,0.75) {$\widehat 2$};
  \begin{scope}[xshift=7.2cm]
  \draw[thick] (-2.2,-1.0) -- (2.2,-1.0);
  \draw[thick] (-2.2,0.8) -- (2.2,0.8);
  \draw[thick, omExo] (-1.6,-1.6) -- (1.6,1.5);
  \draw[omDef, ->] (-0.28,-1.0) ++(0.75,0) arc (0:44:0.75);
  \node[omDef, font=\small] at (0.85,-0.62) {$\widehat a$};
  \draw[omDef, ->] (0.63,0.8) ++(0.75,0) arc (0:44:0.75);
  \node[omDef, font=\small] at (1.76,1.18) {$\widehat b$};
  \end{scope}
\end{tikzpicture}

{\small 左图：对顶角 $\widehat 1 = \widehat 3$
（二者均与 $\widehat 2$ 互补）。右图：两条平行线被
第三条直线所截，同位角 $\widehat a$ 与 $\widehat b$
相等。}
\end{omfigure}

\section{三角形的内角和}

\begin{theorem}[三角形内角和]\label{thm:g7:triangles:sum}
在任意三角形中，三个内角之和为 $180^\circ$。
\end{theorem}

\begin{proof}
设 $ABC$ 为三角形，过 $A$ 作
与 $(BC)$ 平行的直线。在顶点 $A$ 处，三个角沿此平行线排列，
构成一个平角 $180^\circ$：中间的角为
$\widehat{A}$，两侧的角分别与
$\widehat B$、$\widehat C$ 为内错角（平行线！），
因此它们分别等于 $\widehat B$ 与 $\widehat C$。由此可得
$\widehat B + \widehat A + \widehat C = 180^\circ$。
\end{proof}

\begin{omfigure}
\begin{tikzpicture}[scale=1.15]
  \coordinate (B) at (0,0);
  \coordinate (C) at (4.2,0);
  \coordinate (A) at (1.6,2.0);
  \draw[thick] (B) -- (C) -- (A) -- cycle;
  \draw[thick, omExo] (-0.3,2.0) -- (3.9,2.0);
  \node[above, font=\small] at (A) {$A$};
  \node[below left, font=\small] at (B) {$B$};
  \node[below right, font=\small] at (C) {$C$};
  \draw[omDef, ->] (0.55,0) arc (0:51.3:0.55);
  \node[omDef, font=\small] at (0.95,0.32) {$\widehat B$};
  \draw[omMeth, ->] (3.68,0) arc (180:142.4:0.52);
  \node[omMeth, font=\small] at (3.25,0.3) {$\widehat C$};
  \draw[omDef, ->] (A) ++(-0.55,0) arc (180:231.3:0.55);
  \node[omDef, font=\small] at (0.85,1.75) {$\widehat B$};
  \draw[omMeth, ->] (A) ++(0.52,0) arc (0:-37.6:0.52);
  \node[omMeth, font=\small] at (2.35,1.78) {$\widehat C$};
\end{tikzpicture}

{\small 证明图示：沿与 $(BC)$ 平行且过 $A$ 的直线，
角 $\widehat B$、$\widehat A$、$\widehat C$（内错角，颜色对应）
恰好拼成一个平角。}
\end{omfigure}

\begin{example}\label{ex:g7:triangles:sum}
一个三角形的两个角分别为 $67^\circ$ 和 $58^\circ$，则第三个角的
度数为
\[
180 - (67 + 58) = 180 - 125 = 55^\circ .
\]
需要记住的几种特殊情形：等边三角形的每个角
都是 $180 \div 3 = 60^\circ$；直角三角形的两个锐角
互余（和为 $180 - 90 = 90^\circ$）；等腰三角形的底角
相等，因此每个底角
为 $\frac{180 - \text{顶角}}{2}$。
\end{example}

\begin{example}[两个角为直角的三角形不存在]\label{ex:g7:triangles:tworight}
一个三角形若有两个角为 $90^\circ$，则已用去 $180^\circ$，
第三个角只剩 $0^\circ$，这不可能。内角和
这一条件一次就排除了许多三角形。
\end{example}

\section{三角形不等式}

\begin{theorem}[三角形不等式]\label{thm:g7:triangles:inequality}
在任意三角形中，每条边都小于其余两边
之和：设三角形三边为 $a$、$b$、$c$，则
\[
a < b + c .
\]
反之，三个长度中若\emph{最大者}小于其余两者之和，
则总能以这三个长度构成一个三角形。若最大者
\emph{等于}其余两者之和，则所谓“三角形”是退化的：三个
顶点共线。
\end{theorem}
\admitted

\begin{example}\label{ex:g7:triangles:inequality}
三边能否为 $10$、$3$ 和 $2$？检验最大边：$10 < 3 + 2 = 5$ 成立吗？
不成立，这样的三角形不存在。圆规可以说明其原因：分别以线段两端为圆心、半径为
$3$ 和 $2$ 的两段圆弧，从长为 $10$ 的线段两端画出，
永不相交。

三边为 $7$、$5$ 和 $4$ 呢？最大边：$7 < 5 + 4 = 9$，成立，三角形存在
（只需检验最大边即可）。
\end{example}

\begin{omfigure}
\begin{tikzpicture}[scale=0.72]
  \draw[thick] (0,0) -- (10*0.72,0);
  \fill (0,0) circle (1.5pt) node[below, font=\small] {$A$};
  \fill (7.2,0) circle (1.5pt) node[below, font=\small] {$B$};
  \draw[omDef, thick] (2.16,0) arc (0:60:2.16);
  \draw[omDef, thick] (2.16,0) arc (0:-30:2.16);
  \draw[omThm, thick] (7.2,0) ++(-1.44,0) arc (180:120:1.44);
  \draw[omThm, thick] (7.2,0) ++(-1.44,0) arc (180:210:1.44);
  \node[omDef, font=\small] at (1.1,1.6) {半径 $3$};
  \node[omThm, font=\small] at (6.4,1.3) {半径 $2$};
\end{tikzpicture}

{\small 试以 $10$、$3$、$2$ 为边作三角形：两段
圆弧始终无法相交，因为 $3 + 2 < 10$。}
\end{omfigure}

\begin{method}[三角形的存在性与作图]\label{met:g7:triangles:construct}
给定三个长度：
\begin{enumerate}
  \item 将最大者与其余两者之和相比较；若最大者
        不严格小于该和，则停止：不存在三角形；
  \item 否则，按\cref{met:g6:shapes:construct}作图
        （先作底边线段，再作两段圆弧）；
  \item 凡由角作出的三角形，作图后应心算检验
        所给各角之和小于 $180^\circ$——加上第三个角后
        恰为 $180^\circ$。
\end{enumerate}
\end{method}

\section{练习}

\begin{exercise}[$\star$]\label{exo:g7:triangles:1}
求下列各角的余角（和为 $90^\circ$）与补角（和为
$180^\circ$）：$30^\circ$；\ $45^\circ$；\ $72^\circ$。
\end{exercise}

\begin{exercise}[$\star$]\label{exo:g7:triangles:2}
两条直线相交，其中一个角为 $35^\circ$。求
其余三个角的度数，并说明每个角的理由。
\end{exercise}

\begin{exercise}[$\star$]\label{exo:g7:triangles:3}
求下列三角形中第三个角的度数，已知其前两个角的度数分别为：
(a) $40^\circ$ 和 $60^\circ$；\
(b) $90^\circ$ 和 $28^\circ$；\
(c) $75^\circ$ 和 $75^\circ$。
\end{exercise}

\begin{exercise}[$\star$]\label{exo:g7:triangles:4}
一个等腰三角形的顶角为 $40^\circ$，求
其两个底角。另一个等腰三角形的一个底角为
$70^\circ$，求其顶角。
\end{exercise}

\begin{exercise}[$\star$]\label{exo:g7:triangles:5}
下列哪些数组可以作为三角形的三边？请依据
三角形不等式（注意检验最大边）说明理由：
\[
(6,\ 8,\ 12); \qquad
(5,\ 5,\ 11); \qquad
(4,\ 9,\ 13); \qquad
(7,\ 7,\ 7).
\]
\end{exercise}

\begin{exercise}[$\star$]\label{exo:g7:triangles:6}
一个三角形的三个角分别为 $x$、$2x$ 和 $3x$。求 $x$ 及这三个
角的度数。这是什么三角形？
\end{exercise}

\begin{exercise}[$\star$]\label{exo:g7:triangles:7}
作三角形 $ABC$，使 $BC = 6$ cm，
$\widehat{ABC} = 50^\circ$，$\widehat{ACB} = 60^\circ$（用量角器分别在
$B$、$C$ 处画出这两个角）。作图前，先预测
$\widehat{BAC}$ 的度数，再在图上检验。
\end{exercise}

\begin{exercise}[$\star\star$]\label{exo:g7:triangles:8}
两条平行直线被第三条直线所截，形成一个
$54^\circ$ 的角（指在第一个交点处，截线与其中一条
平行线之间的角）。画出图形，并求出
所形成的全部八个角的度数。
\end{exercise}

\begin{exercise}[$\star\star$]\label{exo:g7:triangles:9}
三角形的两边长分别为 $8$ cm 和 $3$ cm，第三边的长度必须介于哪两个
数值之间？给出一个符合条件的整数长度，
以及一个不符合条件的整数长度。
\end{exercise}

\begin{exercise}[$\star\star$]\label{exo:g7:triangles:10}
在三角形 $ABC$ 中，$\widehat A = 2\widehat B$，且
$\widehat C = 90^\circ$。求 $\widehat A$ 和 $\widehat B$。
\end{exercise}

\begin{exercise}[$\star\star\star$]\label{exo:g7:triangles:11}
任意四边形 $ABCD$ 的三个角可以通过
沿一条对角线将其剪成两个三角形来研究。利用此方法证明
每个四边形的四个角之和为 $360^\circ$，并
由此推出五边形（五条边）的内角和。
\end{exercise}

\section{问题：绕街区行走}

\begin{problem}[{周末问题 --- 任意多边形的角：
内角和为 $(n-2) \times 180^\circ$，外部转角之和
始终为 $360^\circ$，以及为何只有三种正多边形能铺满地面}]
\label{pb:g7:triangles:1}
\Cref{exo:g7:triangles:11}中将四边形剪成两个三角形，
并得到 $360^\circ$。本题将这一思路推广到任意边数的多边形，
然后发现另一种完全不同的证明，
即\emph{沿着图形走一圈}，把各处
转过的角度加起来。最后回到你家浴室的地面：在所有
正多边形中，恰有三种可以无缝隙、无重叠地铺满平面，
而蜜蜂已经选定了其中一种。

\medskip
\noindent\textbf{第一部分 --- 内部的角。}
\begin{enumerate}
  \item 从六边形（六条边）的一个顶点出发，画出该顶点引出的所有
        对角线。六边形被剪成了几个三角形？
        它的\emph{全部}内角之和是多少
        （见\cref{thm:g7:triangles:sum}）？
  \item 推广：从边数为 $n$ 的多边形的一个顶点出发，
        这组对角线构成的扇形可将其分成几个三角形？
        由此推出内角和的公式，
        并分别对 $n = 3$、$4$、$5$、$6$ 进行检验。
  \item 计算十边形（$10$ 条边）与
        十二边形（$12$ 条边）的内角和。
  \item 在\emph{正}多边形中，各内角都
        相等。计算等边三角形的内角，
        以及正方形、正五边形、正六边形
        和正八边形的内角。
  \item 不用任何公式，说明多边形
        每增加一个
        顶点，内角和恰好增加 $180^\circ$。
\end{enumerate}

\medskip
\noindent\textbf{第二部分 --- 绕街区行走。}
沿多边形街区的边界行走，始终向前。在
每个拐角处你都要\emph{转过}一个角 --- 即该拐角的\emph{外
角}；转弯持续进行，直到你回到出发点，
面朝出发时的方向。
\begin{enumerate}[resume]
  \item 在内角为 $\widehat A$ 的拐角处，
        转过的角度是多少？（内角与外角互为
        补角。）请计算
        内角分别为 $40^\circ$、
        $60^\circ$、$80^\circ$ 的三角形的三个转角。
  \item 将第 6 题中的三个转角相加，你观察到
        什么关于整周角的事实？
  \item 说明为什么绕一周走遍
        \emph{任意}多边形（三条边或三十条边）时，转角之和总是恰好为
        $360^\circ$。（把所有转角加起来，回到出发点时，
        你的朝向转过了多少？）
  \item 再次推导内角和公式：在
        $n$ 个拐角的每一处，内角 $+$ 转角 $= 180^\circ$；对
        所有拐角求和，并利用第 8 题的结论。
  \item 在正多边形中，所有转角相等，因此每个
        外角为 $\frac{360^\circ}{n}$。利用此结果，
        立即回答：哪种正多边形的内角为
        $150^\circ$？为什么\emph{没有}正多边形
        的内角为 $155^\circ$？
\end{enumerate}

\medskip
\noindent\textbf{第三部分 --- 铺设地面。}
在铺好的地面上，围绕每一点的相接瓷砖的角
之和必须恰为 $360^\circ$ --- 无缝隙，无重叠
（见\cref{pb:g6:lines:1}）。
\begin{enumerate}[resume]
  \item 对等边三角形、正方形和
        正六边形，铺设时在一个角点处有几块相同的瓷砖相接？
        验证每种情形下的角之和都为 $360^\circ$。
  \item 说明正五边形无法铺满地面：三个角
        之和是多少？若是四个角呢？
  \item 说明边数\emph{多于}六条的正多边形
        都无法铺满地面：每个角点处至少有三块瓷砖相接
        （每个内角都小于 $180^\circ$），那么若三个角
        都不小于 $135^\circ$，会发生什么？由此得出
        正多边形铺满地面的完整清单。
  \item 浴室地面常将正八边形与小正方形搭配铺设。
        验证该角点成立：
        $135^\circ + 135^\circ + 90^\circ$。另一种经典搭配是
        在每个角点处混用一个正方形、一个正六边形和一个
        正十二边形：也请验证（十二边形的内角
        可用第 10 题的方法求得）。
  \item 蜜蜂的选择：蜂巢的巢室是正六边形。
        验证其角点（$3$ 个角）的和，并用一句话说明，
        借助狄多的发现
        （见\cref{pb:g6:measure:1}）——为什么在三种可能的
        正多边形中，正六边形盛放同样多蜂蜜时
        所用的蜡最省。
\end{enumerate}
\end{problem}
